\documentclass[11pt]{article}

\usepackage[T1]{fontenc}
\usepackage{lmodern}
\usepackage{amsmath,amssymb,amsthm}
\usepackage[margin=1in]{geometry}
\usepackage{microtype}

\newtheorem{theorem}{Theorem}
\newtheorem{lemma}[theorem]{Lemma}

\newcommand{\R}{\mathbb R}
\newcommand{\Ibar}{\overline{\mathcal I}}
\newcommand{\Bcal}{\mathcal B}
\newcommand{\Rbar}{\overline R}

\title{A purely analytic lower bound for the aggregate reserve at \(K=200\)}
\author{}
\date{}

\begin{document}
\maketitle

\begin{abstract}
This note proves analytically one base-frequency sign in the aggregate-tail
argument.  Every definition is reproduced below.  The proof uses one
concave tangent bound, one Young inequality,
elementary rational enclosures, and two exact Bernstein-basis identities.
In fact, it proves the uniform, hand-checkable estimate
\[
 \mathcal B(\rho,200)>60
 \qquad\left(\frac7{51}\leq\rho\leq\frac{39}{50}\right).
\]
No decimal approximation or machine-generated sign enclosure is a premise.
\end{abstract}

\section{The exact assertion}

For \(0<\rho<1\), put
\[
 G_1(s)=\frac{\sqrt{1-s^2}-s\arccos s}{\pi},
 \qquad
 \eta_\rho=G_1\!\left(\max\{\rho,\tfrac12\}\right),
 \qquad
 d_\rho=\frac{2\arcsin\rho}{\pi}.
\]
At the fixed frequency \(K=200\), define
\[
 \mu=200\rho,
 \qquad
 b=\frac{2\pi\rho}{1-\rho},
 \qquad
 c=bK=\frac{400\pi\rho}{1-\rho},
 \qquad
 \Rbar=\mu+\frac12,
 \qquad
 S=\sqrt{\Rbar^2+c},
\]
and
\begin{equation}
 \Ibar
 =\frac12\left[
      \Rbar S+c\operatorname{arsinh}\frac{\Rbar}{\sqrt c}
      -\Rbar^2
    \right].
 \label{eq:I-definition}
\end{equation}
These are exactly the specializations at \(K=200\) of the aggregate-tail
quantities in the main proof.  The reserve to be estimated is
\begin{equation}
\begin{aligned}
 \Bcal(\rho,200)={}&
 (\mu-\tfrac12)(200\eta_\rho-1)
 +\frac{d_\rho}{2}(\mu-\tfrac32)(\mu-\tfrac12)\\
 &-(1+d_\rho)\Ibar
 -\frac{8(\mu+\tfrac12)}{15\sqrt\mu}.
\end{aligned}
\label{eq:B-definition}
\end{equation}

\begin{theorem}
\label{thm:B-positive}
For every
\[
       \frac7{51}\leq\rho\leq\frac{39}{50},
\]
one has
\[
       \Bcal(\rho,200)>60>0.
\]
\end{theorem}

\section{A concave majorant for the integral term}

Set
\[
 H(u)=\frac12\left[
 u\sqrt{1+u^2}+\operatorname{arsinh}u-u^2
 \right]
 =\int_0^u\bigl(\sqrt{1+t^2}-t\bigr)\,dt.
\]
Scaling \eqref{eq:I-definition} by
\(u=\Rbar/\sqrt c\) gives the exact identity
\begin{equation}
       \Ibar=cH(u).
       \label{eq:I-scaled}
\end{equation}

\begin{lemma}[Concave tangent]
\label{lem:H-tangent}
For every \(u>0\),
\begin{equation}
       H(u)<\frac6{25}u+\frac{12}{25}.
       \label{eq:H-tangent}
\end{equation}
\end{lemma}

\begin{proof}
One has
\[
 H'(u)=\sqrt{1+u^2}-u,
 \qquad
 H''(u)=\frac{u}{\sqrt{1+u^2}}-1<0.
\]
Thus every tangent line is a global upper bound.  Choose
\[
 u_0=\frac{589}{300},
 \qquad
 \sqrt{1+u_0^2}=\frac{661}{300}.
\]
Then
\[
 H'(u_0)=\frac6{25},
 \qquad
 \operatorname{arsinh}u_0
 =\log\left(u_0+\sqrt{1+u_0^2}\right)
 =\log\frac{25}{6},
\]
and hence
\begin{equation}
 H(u_0)-\frac6{25}u_0
 =\frac12\left(\log\frac{25}{6}-\frac{589}{1250}\right).
 \label{eq:tangent-intercept}
\end{equation}

For completeness, the elementary series
\[
 \log2=2\sum_{n=0}^{\infty}
       \frac{1}{(2n+1)3^{2n+1}}
\]
gives
\[
 \log2
 <\frac23+\frac23\sum_{n=1}^{\infty}3^{-(2n+1)}
 =\frac{25}{36}.
\]
Also \(\log(1+x)<x\) for \(x>0\).  Therefore
\[
 \log\frac{25}{6}
 =2\log2+\log\frac{25}{24}
 <\frac{25}{18}+\frac1{24}
 =\frac{103}{72}
 <\frac{1789}{1250};
\]
the last difference is \(29/45000\).  Substitution in
\eqref{eq:tangent-intercept} makes the tangent intercept strictly smaller
than \(12/25\), proving \eqref{eq:H-tangent} by concavity.
\end{proof}

We shall rationalize the remaining square root using the elementary Young
inequality
\begin{equation}
       \sqrt c\leq\frac{c}{2a}+\frac a2
       \qquad(a>0),
       \label{eq:young}
\end{equation}
which is equivalent to \((\sqrt c-a)^2\geq0\).  The classical bound
\begin{equation}
       \pi<\frac{22}{7}
       \label{eq:pi-upper}
\end{equation}
follows, for example, from
\[
 0<\int_0^1\frac{x^4(1-x)^4}{1+x^2}\,dx=\frac{22}{7}-\pi.
\]
Consequently
\[
       c<c_+(\rho):=\frac{8800}{7}\frac{\rho}{1-\rho}.
\]
Combining Lemma~\ref{lem:H-tangent} with \eqref{eq:young} gives, for every
\(a>0\),
\begin{equation}
 \Ibar<Q_a(\rho):=
 \frac{3\Rbar c_+(\rho)}{25a}
 +\frac{3a\Rbar}{25}
 +\frac{12c_+(\rho)}{25}.
 \label{eq:Qa}
\end{equation}
All quantities on the right are rational functions of \(\rho\) when
\(a\) is rational.

\section{The coefficient of \(d_\rho\) is positive}

The replacement of \(d_\rho\) by an elementary lower bound requires the
sign of its coefficient.  We verify that sign before making any such
replacement.

\begin{lemma}
\label{lem:d-coefficient}
On \([7/51,39/50]\),
\begin{equation}
 \frac12(\mu-\tfrac32)(\mu-\tfrac12)-\Ibar>0.
 \label{eq:d-coefficient-positive}
\end{equation}
\end{lemma}

\begin{proof}
Write
\[
 \Phi(\rho)=\Rbar^2\frac{1-\rho}{\rho}.
\]
Since
\[
 \frac{d}{d\rho}\log\Phi(\rho)
 =\frac{200\rho-400\rho^2-1/2}
        {\rho(1-\rho)\Rbar},
\]
its two critical points are
\[
       q_\pm=\frac{10\pm7\sqrt2}{40}.
\]
The elementary bounds \(7/5<\sqrt2<3/2\) show
\[
 q_-<\frac1{200}<\frac7{51}<q_+<\frac{41}{80}<\frac{39}{50}.
\]
Thus \(\Phi\) increases and then decreases on the interval in question,
so its minimum is attained at one of the two endpoints.  From
\eqref{eq:pi-upper},
\[
 u^2=\frac{\Rbar^2(1-\rho)}{400\pi\rho}
 >\frac{7\Rbar^2(1-\rho)}{8800\rho}.
\]
At the two endpoints the rational lower bounds on the right are
\[
 \frac{8128201}{2080800},
 \qquad
 \frac{685783}{124800},
\]
respectively, and
\[
 \frac{8128201}{2080800}-\frac{225}{64}
 =\frac{1625777}{4161600}>0,
 \qquad
 \frac{685783}{124800}-\frac{225}{64}
 =\frac{247033}{124800}>0.
\]
Hence \(u>15/8\) throughout the interval.

By \eqref{eq:I-scaled} and Lemma~\ref{lem:H-tangent},
\[
 \frac{\Ibar}{\Rbar^2}
 =\frac{H(u)}{u^2}
 <\frac6{25u}+\frac{12}{25u^2}
 \leq\frac{496}{1875}.
\]
Moreover \(\Rbar\geq2851/102>27\), and the function
\[
 \frac{(R-2)(R-1)}{2R^2}
 =\frac12-\frac{3}{2R}+\frac1{R^2}
\]
is increasing for \(R>4/3\).  Therefore
\[
 \frac{(\mu-\tfrac32)(\mu-\tfrac12)}{2\Rbar^2}
 \geq\frac{25\cdot26}{2\cdot27^2}
 =\frac{325}{729}.
\]
Finally
\[
       \frac{325}{729}-\frac{496}{1875}
       =\frac{82597}{455625}>0,
\]
which proves \eqref{eq:d-coefficient-positive}.
\end{proof}

\section{Rational lower envelopes}

We use only the consequences
\[
 \frac1\pi>\frac7{22},
 \qquad
 \sqrt2>\frac75,
 \qquad
 \sqrt3>\frac53
\]
of elementary positive-side squaring and \eqref{eq:pi-upper}.

On \(7/51\leq\rho\leq1/2\), one has
\[
 \eta_\rho=G_1(1/2)
 =\frac{\sqrt3}{2\pi}-\frac16
 >\frac{13}{132},
 \qquad
 d_\rho=\frac{2\arcsin\rho}{\pi}>\frac7{11}\rho.
\]
Also \(1400/51\leq\mu\leq100\), and
\((26/5)^2<1400/51\), so
\begin{equation}
 \frac{8(\mu+1/2)}{15\sqrt\mu}
 =\frac8{15}\left(\sqrt\mu+\frac1{2\sqrt\mu}\right)
 <\frac8{15}\left(10+\frac5{52}\right)
 =\frac{70}{13}.
 \label{eq:E-low}
\end{equation}
Define
\begin{equation}
\begin{aligned}
 \eta_<&=\frac{13}{132},
 &d_<&=\frac7{11}\rho,
 &E_<&=\frac{70}{13},\\
 L_< (\rho)
 &=(\mu-\tfrac12)(200\eta_<-1)
 +\frac{d_<}{2}(\mu-\tfrac32)(\mu-\tfrac12)
 -(1+d_<)Q_{18}(\rho)-E_<.
\end{aligned}
\label{eq:L-low}
\end{equation}

On \(1/2\leq\rho\leq39/50\), the integral representation
\[
 \eta_\rho=\frac1\pi\int_\rho^1\arccos t\,dt
\]
and \(\arccos t\geq\sqrt{2(1-t)}\) give
\[
 \eta_\rho\geq\frac{2\sqrt2}{3\pi}(1-\rho)^{3/2}.
\]
Here \(1-\rho\geq11/50>(23/50)^2\), and consequently
\begin{equation}
 \eta_\rho>\frac{1127}{8250}(1-\rho)=:\eta_>.
 \label{eq:eta-high}
\end{equation}
The Maclaurin series for \(\arcsin\rho\) has positive coefficients, so
\begin{equation}
 d_\rho>\frac7{11}P_7(\rho)=:d_>,
 \qquad
 P_7(\rho)=\rho+\frac{\rho^3}{6}
 +\frac{3\rho^5}{40}+\frac{5\rho^7}{112}.
 \label{eq:d-high}
\end{equation}
Finally \(100\leq\mu\leq156<169\), whence
\begin{equation}
 \frac8{15}\left(\sqrt\mu+\frac1{2\sqrt\mu}\right)
 <\frac8{15}\left(13+\frac1{20}\right)
 =\frac{174}{25}<7.
 \label{eq:E-high}
\end{equation}
Define
\begin{equation}
\begin{aligned}
 E_>&=7,\\
 L_>(\rho)
 &=(\mu-\tfrac12)(200\eta_>-1)
 +\frac{d_>}{2}(\mu-\tfrac32)(\mu-\tfrac12)
 -(1+d_>)Q_{67}(\rho)-E_>.
\end{aligned}
\label{eq:L-high}
\end{equation}

Lemma~\ref{lem:d-coefficient}, the lower bounds for \(d_\rho\) and
\(\eta_\rho\), the upper bounds for the last term, and
\eqref{eq:Qa} now give the logically ordered implications
\begin{equation}
 \Bcal(\rho,200)>L_<(\rho)
 \quad\left(\frac7{51}\leq\rho\leq\frac12\right),
 \qquad
 \Bcal(\rho,200)>L_>(\rho)
 \quad\left(\frac12\leq\rho\leq\frac{39}{50}\right).
 \label{eq:B-L}
\end{equation}
The overlap at \(\rho=1/2\) is harmless.

\section{Exact Bernstein closure}

For \(0\leq x\leq1\), write
\[
 B_{k,n}(x)=\binom nkx^k(1-x)^{n-k}.
\]
These functions are nonnegative and sum to one.  The conversion from a
power polynomial \(\sum_{j=0}^nq_jx^j\) to Bernstein coefficients is the
finite rational identity
\begin{equation}
 \sum_{j=0}^nq_jx^j
 =\sum_{k=0}^n\left(
   \sum_{j=0}^k\frac{\binom kj}{\binom nj}q_j
  \right)B_{k,n}(x).
 \label{eq:bernstein-conversion}
\end{equation}
Thus every entry below can be checked by additions and cross
multiplications of integers.

For the lower interval set
\[
 x_<=\frac{\rho-7/51}{1/2-7/51}=\frac{102\rho-14}{37}.
\]
Direct substitution of \eqref{eq:Qa} and \eqref{eq:L-low}, followed by
\eqref{eq:bernstein-conversion}, gives the exact identity
\begin{equation}
 (1-\rho)L_<(\rho)
 =\frac1{D_<}\sum_{k=0}^{4}N_{<,k}B_{k,4}(x_<),
 \qquad
 D_<=10835145921600,
 \label{eq:low-bernstein}
\end{equation}
where
\[
\begin{array}{c|r}
k&N_{<,k}\\ \hline
0&2977383044192480\\
1&4628920786869499\\
2&5747338995812809\\
3&6988840660610010\\
4&6621947656848702
\end{array}
\]
The smallest entry is \(N_{<,0}\), and
\[
 N_{<,0}-60D_<=2327274288896480>0.
\]
It follows from \eqref{eq:low-bernstein} that
\begin{equation}
       (1-\rho)L_<(\rho)>60.
       \label{eq:low-sixty}
\end{equation}

For the upper interval set
\[
 x_>=\frac{\rho-1/2}{39/50-1/2}
 =\frac{25}{7}\left(\rho-\frac12\right).
\]
The corresponding exact identity obtained from \eqref{eq:Qa},
\eqref{eq:L-high}, and \eqref{eq:bernstein-conversion} is
\begin{equation}
 (1-\rho)L_>(\rho)
 =\frac1{D_>}\sum_{k=0}^{10}N_{>,k}B_{k,10}(x_>),
 \qquad
 D_>=2487375000000000000000,
 \label{eq:high-bernstein}
\end{equation}
where
\[
\begin{array}{c|r}
k&N_{>,k}\\ \hline
0&873011369240936279296875\\
1&881095006796855712890625\\
2&884923416376472412109375\\
3&881704924277462068359375\\
4&867541422038412498046875\\
5&837327296894843418528125\\
6&784608885866591053548975\\
7&701387295930232136634075\\
8&577837649792730661053099\\
9&401902155200474706698937\\
10&158688669397963044788415
\end{array}
\]
The smallest entry is \(N_{>,10}\), and
\[
 N_{>,10}-60D_>
 =9446169397963044788415>0.
\]
Therefore
\begin{equation}
       (1-\rho)L_>(\rho)>60.
       \label{eq:high-sixty}
\end{equation}

Since \(0<1-\rho<1\), both \eqref{eq:low-sixty} and
\eqref{eq:high-sixty} imply that the corresponding \(L\)-function is
strictly larger than \(60\).  Combining this conclusion with
\eqref{eq:B-L} proves Theorem~\ref{thm:B-positive}.

\paragraph{Verification boundary.}
The two Bernstein identities are finite polynomial identities over
\(\mathbb Q\).  Formula \eqref{eq:bernstein-conversion} is a direct
hand-checking recipe for every printed integer.  No interval evaluation,
floating-point comparison, or unprinted case is used in the proof.

\end{document}
